\subsection{Data mart và các thao tác OLAP}
Trên kho dữ liệu hình sao, nhóm dựng bốn data mart theo chủ đề: bán hàng, xử lý đơn, trải nghiệm khách và thu hút seller. Mỗi mart là một lớp view trong schema riêng, tổng hợp sẵn các độ đo mà dashboard cần (GMV theo ngày, danh mục, bang; thời gian từng chặng và tỷ lệ giao trễ; điểm đánh giá theo mức trễ, cơ cấu thanh toán, RFM, cohort; funnel từ khách tiềm năng đến seller có đơn và scorecard seller). Các view có chi phí tính toán lớn được cài dưới dạng materialized view và được làm mới trong mỗi lần chạy pipeline, nên mọi truy vấn trên mart đều trả lời dưới 0,2 giây. Các thao tác OLAP được minh hoạ bằng SQL theo mô hình ROLAP: roll-up bằng ROLLUP, GROUPING SETS và CUBE; drill-down bằng cách thêm cấp thấp hơn của phân cấp vào GROUP BY; slice bằng cách cố định một chiều (năm 2017); dice bằng cách giới hạn đồng thời nhiều chiều (vùng Sudeste, nhóm ngành Home \& Furniture, quý 4); pivot bằng tổng có điều kiện FILTER. Ví dụ, lát cắt năm 2017 cho thấy vùng Sudeste chiếm 64,3\% GMV, và drill-down tháng 11/2017 chỉ ra ngày Black Friday 24/11 có GMV cao nhất tháng (152.653,74 BRL). Nhóm chọn ROLAP thay vì MOLAP vì hệ thống dùng PostgreSQL, dữ liệu vài trăm nghìn dòng fact, kho luôn cập nhật sau mỗi lần nạp tăng dần mà không cần xử lý lại khối đa chiều, còn Power BI đảm nhiệm lớp ngữ nghĩa và tương tác kéo thả cho người dùng cuối. Lớp ngữ nghĩa đó gồm 32 quan hệ, các phân cấp thời gian, địa lý, sản phẩm, seller và thư viện measure DAX cho các KPI chính, được đối chiếu với giá trị tính bằng SQL.

{\small\begin{longtable}{llllp{6.3cm}}
\caption{Các data mart: view, loại, số dòng và mô tả}\label{tab:marts}\\
\toprule
\textbf{Mart} & \textbf{View} & \textbf{Loại} & \textbf{Dòng} & \textbf{Mô tả} \\
\midrule
\endfirsthead
\toprule
\textbf{Mart} & \textbf{View} & \textbf{Loại} & \textbf{Dòng} & \textbf{Mô tả} \\
\midrule
\endhead
\bottomrule
\endfoot
mart\_acquisition & v\_funnel\_by\_origin\_month & view & 113 & Funnel theo tháng liên hệ đầu x origin: MQL, đã chốt, đã có đơn, tỷ lệ, số ngày TB, GMV 90 ngày đầu \\
mart\_acquisition & v\_seller\_by\_channel & materialized & 7 & Theo kênh thu hút seller: số seller, số seller có đơn, GMV, GMV / seller, trễ, review; cờ mẫu nhỏ (< 30 seller có đơn) \\
mart\_acquisition & v\_seller\_scorecard & materialized & 3.095 & 1 dòng / seller có đơn: GMV, đơn, trễ, hủy, review, hạng và tỷ trọng luỹ kế (Pareto), cờ mẫu nhỏ (< 10 đơn) \\
mart\_cx & v\_customer\_cohort & materialized & 220 & Cohort theo tháng mua đầu tiên: số khách còn mua lại sau 0,1,2... tháng và tỷ lệ quay lại \\
mart\_cx & v\_customer\_rfm & materialized & 94.983 & 1 dòng / khách thật (customer\_unique\_id): recency, frequency, monetary, điểm R-F-M 1-5 và phân khúc; đơn không thuộc nhóm Failed \\
mart\_cx & v\_payment\_mix\_month & view & 90 & Cơ cấu thanh toán theo tháng mua x hình thức: số lần, giá trị, trả góp TB, tỷ trọng giá trị \\
mart\_cx & v\_review\_vs\_delay & view & 105 & Điểm review và \% tiêu cực theo nhóm trễ giao (giao sớm ... trễ trên 7 ngày) x tháng mua; chỉ đơn đã giao \\
mart\_fulfillment & v\_fulfillment\_order & view & 99.441 & 1 dòng / đơn: mốc ngày, thời gian từng chặng, số ngày hứa và trễ, nhóm trễ (delay\_bucket), bang và vùng của khách \\
mart\_fulfillment & v\_fulfillment\_state\_month & view & 565 & Theo tháng mua x bang: số đơn, tỷ lệ đúng hạn / hủy, thời gian trung bình từng chặng (chỉ đơn đã giao) \\
mart\_fulfillment & v\_shipping\_distance & view & 112.650 & Khoảng cách seller -> khách (km, haversine) cho từng dòng hàng, nhóm 0-100 / 100-500 / 500-1000 / trên 1000 km; distance\_km NULL nếu thiếu toạ độ (has\_distance = false) \\
mart\_fulfillment & v\_status\_funnel & view & 125 & Số đơn đi qua từng chặng (đặt, duyệt, bàn giao, giao, hủy) theo tháng mua, dạng dài \\
mart\_sales & v\_category\_pareto & view & 74 & Pareto 80/20 của GMV theo danh mục: xếp hạng, tỷ trọng, tỷ trọng luỹ kế, hạng A/B/C \\
mart\_sales & v\_geo\_state & view & 27 & Dữ liệu bản đồ theo bang: toạ độ tâm (trung bình các zip có toạ độ), số khách, đơn, GMV, freight ratio, số seller, lead time, tỷ lệ trễ \\
mart\_sales & v\_geo\_zip & view & 14.889 & Dữ liệu bản đồ theo zip có toạ độ và có khách hoặc seller: toạ độ, thành phố, số đơn và GMV phía khách, số seller \\
mart\_sales & v\_sales\_category\_month & view & 1.282 & GMV, đơn, dòng hàng theo tháng x nhóm ngành hàng x danh mục (số đơn không cộng được giữa các danh mục) \\
mart\_sales & v\_sales\_daily & view & 614 & GMV, số đơn, số dòng hàng, phí vận chuyển, AOV theo ngày mua (chỉ đơn không thuộc nhóm Failed) \\
mart\_sales & v\_sales\_state\_month & view & 558 & GMV, đơn, phí vận chuyển, tỷ lệ freight/price theo tháng x bang của khách \\
\end{longtable}}

\paragraph{Roll-up và pivot trên dữ liệu thật.} Bảng \ref{tab:olap-rollup} là kết quả roll-up bằng \code{ROLLUP} (SQL ở Phụ lục B); Bảng \ref{tab:olap-pivot} là pivot GMV theo vùng và năm bằng tổng có điều kiện \code{FILTER}.
\begin{table}[htbp]
\centering\small
\caption{Roll-up bằng ROLLUP: GMV theo quý, năm và tổng}
\label{tab:olap-rollup}
\begin{tabular}{llrr}
\toprule
\textbf{Năm} & \textbf{Quý} & \textbf{Số đơn} & \textbf{GMV (BRL)} \\
\midrule
2017 & 2017-Q1 & 5.122 & 733.398,94 \\
2017 & 2017-Q2 & 9.222 & 1.286.918,78 \\
2017 & 2017-Q3 & 12.445 & 1.681.949,00 \\
2017 & 2017-Q4 & 17.586 & 2.406.225,55 \\
2017 & (cả năm) & 44.375 & 6.108.492,27 \\
2018 & 2018-Q1 & 20.979 & 2.764.402,78 \\
2018 & 2018-Q2 & 19.897 & 2.849.730,26 \\
2018 & 2018-Q3 & 12.655 & 1.727.049,37 \\
2018 & (cả năm) & 53.531 & 7.341.182,41 \\
Tổng & (cả năm) & 97.906 & 13.449.674,68 \\
\bottomrule
\end{tabular}
\end{table}

\begin{table}[htbp]
\centering\small
\caption{Pivot bằng FILTER: GMV (BRL) theo vùng × năm}
\label{tab:olap-pivot}
\begin{tabular}{lrrrr}
\toprule
\textbf{Vùng} & \textbf{2016} & \textbf{2017} & \textbf{2018} & \textbf{Tổng} \\
\midrule
Sudeste & 27.701 & 3.928.343 & 4.866.293 & 8.822.336 \\
Sul & 7.456 & 883.020 & 1.047.636 & 1.938.111 \\
Nordeste & 6.013 & 727.733 & 804.568 & 1.538.314 \\
Centro-Oeste & 2.356 & 401.380 & 458.090 & 861.825 \\
Norte & 1.200 & 168.018 & 164.596 & 333.814 \\
\bottomrule
\end{tabular}
\end{table}

\todo{Chọn thêm các truy vấn drill-down, slice, dice từ handoff/olap-examples.md; viết đoạn so sánh ROLAP và MOLAP trong bối cảnh bài toán (liên hệ mục 1.6).}
